% 第 24 章　麦克斯韦把电和磁拼成一个整体
\chapter{麦克斯韦把电和磁拼成一个整体}

\step{反常识开场}

晚上回家，你按下开关，台灯亮了。问一个朴素得近乎多余的问题：能量从插座——追根溯源是几百公里外的发电厂——来到灯丝，走的是哪条路？

几乎所有人都会回答：沿导线。电流在导线里流，就像水在水管里流，能量当然也在导线里面跑。这个答案自然得不容置疑：你摸得到导线，看得见它连着开关和灯泡，剪断它灯就灭。

再看一件同样熟悉的事。夏天的正午，阳光晒在皮肤上发烫。太阳和地球之间隔着一亿五千万公里几乎空无一物的空间，没有导线，没有任何物质管道，能量却源源不绝地送过来：每秒送到整个地球的太阳能，折合成功率超过人类全部发电站总功率的几千倍。收音机和手机也一样：天线没有一根线连到电台，节目照播不误。

把两件事放在一起，疑问就立起来了：在没有导线的地方，能量凭什么能穿过真空？而在有导线的地方，能量真的在导线里面走吗？

本章会给你一个让许多人坐不住的答案：就算在最普通的台灯电路里，能量也\textbf{不是}在导线内部流动，而是从导线周围的空间里流进灯泡的——导线更像铁轨，负责规定路线，不负责“装运”能量。把这个答案撑起来的，是麦克斯韦在十九世纪六十年代完成的拼接：电和磁不是两门沾亲带故的学科，而是同一个对象——\textbf{电磁场}——的两个侧面。

\step{先猜一猜}

在继续读之前，请写下你对三个问题的猜测。读完本章我们会回来对答案。

问题一：台灯电路里，能量从电源到灯丝走哪条路？（a）电子像搬运工，在导线里把能量一份份背过去；（b）能量以某种波的形式沿导线内部冲过去；（c）能量在导线外面的空间里走，导线只负责引导方向。

问题二：电现象和磁现象是什么关系？（a）两种独立现象，只在某些场合互相影响；（b）同一种东西的两个侧面。

问题三：变化的电场能不能像电流一样“制造”磁场？在 1860 年以前，所有已知实验里，磁场的来源只有磁铁和电流两类。你的直觉会选“能”还是“不能”？注意：这个问题在当年没有任何直接实验证据，完全靠推理回答。

说说历史处境。1820 年奥斯特发现电流让磁针偏转，1831 年法拉第发现变化的磁场产生电流，此后三十年，电学和磁学各自长成枝繁叶茂的两本书：库仑定律、安培环路定律、法拉第电磁感应定律，每一条都好用。麦克斯韦 1860 年代动手时的全部材料，就是这几条各自成立的定律，外加一个多数人根本没注意到的小裂缝。请记住你自己选的答案——尤其是问题三。

\step{动手实验}

\begin{experiment}[用收音机听电的“喷嚏”]
材料：一台能收中波（AM）的收音机（老式收音机、带收音功能的音箱都行），一只手电筒或带开关的小电扇；如果有一个压电式电子打火机，效果更好。

第一步：把收音机调到中波段没有电台的频率，音量开大，应该听到轻微的沙沙声。

第二步：把收音机放在桌上，在半米之内反复开、关手电筒（或电扇）。仔细听：每次开关的\textbf{瞬间}，收音机里会“咔”的一声。通电期间反而安静。

第三步：换成电子打火机，在收音机旁按动点火，“咔哒”声会明显得多——那是电火花在“喊话”。

第四步：慢慢把打火机拉远，听声音怎样变化；再换到调频（FM）波段试一次，声音会弱得多，甚至完全听不到。

安全提示：只用手边的小功率物件，不要拆插座，不要用市电制造火花。
\end{experiment}

这个实验的奇怪之处值得停下来想清楚：打火机和收音机之间\textbf{没有任何导线连接}。开关闭合那一下，电流从零突然变到某个值，这件事竟然“隔空”被收音机知道了。报信的载体穿过空气——而且从实验上看，它似乎什么介质都不需要。

类似的日常经验其实不少：手机来电前几秒，旁边的音箱常会“滋滋”作响；老式日光灯管拿到高压输电线下方，夜里会微微发亮——灯管两头什么都没接，是空间里的电场在推灯管里的电荷；微波炉加热时，同在 \(2.4\) 吉赫兹频段工作的 WiFi 常会明显变慢——炉门网眼拦不住漏出来的那点电磁扰动，而路由器听得见。这些现象共用同一个前提：空间本身可以携带电的、磁的影响，并且这种携带不需要电线。

请先记下三个观察结果，后面逐个接住：咔哒声出现在开关的\textbf{瞬间}而非通电期间；拉远后声音变弱；FM 波段几乎听不到。

\step{旧模型遇到麻烦}

先清点手里的旧模型。到本章为止，电磁学有四块基石：

电荷产生电场（第 19 章），高斯定律把它改写成总账：从一个闭合面“穿出”多少电场，由面内装了多少净电荷决定；磁场没有“磁荷”作为源头，磁感线总是闭合的（第 22 章）；电流产生环绕导线的磁场，安培环路定律给出总账：沿闭合回路累计一圈磁场，正比于穿过回路所围曲面的电流；变化的磁场产生旋涡状的电场，法拉第定律给出又一笔总账（第 23 章）。

四条定律各自久经考验。麻烦出在把它们拼在一起的时候。

\textbf{麻烦一：一只电容器戳破了安培环路定律。}看充电中的电容器：导线里有电流 \(I\) 流向极板。在导线周围画一个圆圈回路 \(L\)，安培环路定律说：沿 \(L\) 累计的磁场，等于 \(\mu_0\) 乘以穿过“以 \(L\) 为边界的\textbf{任意}曲面”的电流。取一张平平的圆面，导线穿过它，穿过的电流是 \(I\)，没问题。可定律说的是“任意”——那就取一个像碗一样向右鼓起、从电容器两极板之间穿过去的面：极板之间是绝缘的，\textbf{没有任何电荷穿过这个碗面}，穿过的电流是零。同一个回路，同一段磁场，两张面给出 \(I\) 和 \(0\) 两个互相矛盾的答案。定律在自己家里打架。

\begin{center}
\begin{tikzpicture}[>=Stealth]
  \draw[thick] (-0.5,0) -- (2.6,0);
  \draw[thick] (3.4,0) -- (6.3,0);
  \draw[very thick] (2.6,-0.9) -- (2.6,0.9);
  \draw[very thick] (3.4,-0.9) -- (3.4,0.9);
  \node at (3.0,-1.35) {极板间隙};
  \draw[->,thick] (0.2,0.32) -- (1.3,0.32) node[midway,above]{$I$};
  \fill[blue!10] (1.8,0) ellipse (0.5 and 1.3);
  \draw[thick] (1.8,0) ellipse (0.5 and 1.3);
  \node at (1.8,1.7) {回路 $L$};
  \node at (1.8,-1.8) {面 A：导线穿过};
  \draw[dashed,thick] (1.8,1.3) .. controls (3.2,1.7) and (4.4,0.9) .. (4.4,0)
                      .. controls (4.4,-0.9) and (3.2,-1.7) .. (1.8,-1.3);
  \node at (5.0,1.1) {面 B：从极板间鼓过去};
\end{tikzpicture}
\end{center}

这不是吹毛求疵。安培环路定律是从\textbf{稳恒}电流的实验里总结出来的，稳恒电流的回路永远闭合，永远遇不到“断口”。充电电路里电流在极板处断了头，把旧定律外推到这里，就撞了墙。这正是一条经验定律的典型死法：在它被检验过的范围内毫发无损，一旦越界就自相矛盾。

\textbf{麻烦二：对称性的缺口。}法拉第说，变化的磁场能生出电场。自然要问：变化的电场能不能生出磁场？请注意，对称性只是线索，不是证据——宇宙没有义务长得对称。但这个问题和麻烦一隐隐指向同一个补丁。

\textbf{麻烦三：“立即”的幽灵。}在库仑定律里，电荷动一下，远处的力似乎立刻跟着变。第 19 章引入电场，本来就是为了赶走这种超距作用；可旧的定律集合里没有一条说清楚：场的变化怎样传播、以多快的速度传播。旧模型里藏着一封没拆的信。

\step{发明新概念}

麦克斯韦的补丁，来自对那个“碗面”的第二次注视。碗面从两极板之间穿过，确实没有电荷穿过它；但有一样东西在变：\textbf{电场}。导线每送一库仑电荷到上极板，极板间的电场就增强一分。电场增强的速率和导线电流严格成正比——用第 19 章的高斯定律就能看出来：极板间的电通量 \(\varPhi_E\) 等于极板电荷除以 \(\varepsilon_0\)，所以

\[
\varepsilon_0\frac{d\varPhi_E}{dt}=\frac{dQ}{dt}=I,
\]

它正好等于导线里的电流！不多也不少。

这正是位移电流成为关键补丁的原因：它不是为了让方程更漂亮而添加的装饰，而是被两笔旧账——高斯定律和电荷守恒——同时逼出来的。缺了它，安培环路定律在断口处自相矛盾；添了它，所有曲面的账自动对齐。一个补丁同时修好两个故障，这在物理学里通常是“走对了路”的信号。

麦克斯韦提议：在安培环路定律里，把“穿过曲面的电流”换成“传导电流加上 \(\varepsilon_0\) 乘以电通量的变化率”。后一项他起名叫\textbf{位移电流}。这个名字起得并不好（“模型的边界”一节会收回它），但补丁的效果立竿见影：在平面 A 上，传导电流是 \(I\)，电场几乎为零，位移电流项为零，合计 \(I\)；在碗面 B 上，传导电流是零，位移电流恰好是 \(I\)，合计还是 \(I\)。矛盾消失了。而且顺带修好的还有电荷守恒的账：电流在极板“断头”的地方，恰好由变化的电场“接班”，账本的每一页都对得上。

补上这一笔之后，麦克斯韦把当时全部电磁实验定律拼成了四条。先用“区域总账”的语言各说一句，数学式子留到“数学翻译器”：

一、\textbf{电场的高斯定律}：从任何闭合面穿出的电场总量，等于面内净电荷除以 \(\varepsilon_0\)。电荷是电场唯一的源头；电场线起于正电荷、止于负电荷。

二、\textbf{磁场的高斯定律}：从任何闭合面穿出的磁场总量，恒等于零。至今没有发现孤立的磁荷；磁感线没有端点，只能自行闭合或伸向无穷远。

三、\textbf{法拉第定律}：沿任何闭合回路累计一圈的电场（环流），等于回路所围曲面上磁通量变化率的负值。变化的磁场制造旋涡状的电场；负号是楞次定律：感应的效果总是“反对”变化本身。

四、\textbf{安培—麦克斯韦定律}：沿任何闭合回路累计一圈的磁场，等于 \(\mu_0\) 乘以穿过该曲面的“传导电流加位移电流”。电流和变化的电场都能制造旋涡状的磁场。

这四条拼在一起，电和磁就不再分家了。哪怕某个区域里没有电荷、没有电流，第三、四条仍然把电场和磁场绑在一起：一处电场随时间变化，周围必然出现环绕的磁场；一处磁场随时间变化，周围必然出现环绕的电场。这里要特别小心一种流行的简化说法——“电场产生磁场、磁场再产生电场、如此接力前进”。这不是麦克斯韦方程说的。方程说的是：\textbf{场的空间分布和它的时间变化必须同时满足这组约束}，满足约束的场才是合法的。这组约束有一个惊人的后果：它允许一种不依赖任何电荷和电流、能在真空中自我维持、并以确定速度向外传播的场分布。这就是电磁波，第 25 章专门讲它；本章只负责把这块路牌立起来。

场的地位也随之升级。在第 19 章，你还可以把电场当成一种记账工具；现在不行了。电容器极板间空无一物，那里却有真实的物理过程在发生——变化的电场卷起磁场，和电流卷起磁场一模一样；能量可以穿过连空气都没有的真空。场携带能量，也携带动量，它和带电小球一样是物理世界的正式成员，不是草稿纸上的辅助线。

最后交代“总账”这个比喻的边界。它\textbf{像的部分}：四条方程每一条都是“边界上累计的东西，由内部的源头决定”，和水箱的进出账一样，边界一划定，账就跑不掉。它\textbf{不像的部分}：水箱里流的是水，这里边界上累计的是场；内部的“源头”可以是电荷和电流，也可以是另一种场的时间变化；而且这笔账是逐点、逐时刻严格成立的，不是月底结算的近似。

还要补一句后话。麦克斯韦 1865 年发表这套方程时，位移电流和电磁波都只是纸面结论，许多同行半信半疑；直到 1887 年，赫兹才在实验室里用火花装置真正发出并接收到了电磁波——那是“先猜一猜”里问题三的迟到二十多年的实验回答。赫兹实验的细节属于第 25 章，这里只记住一点：本章的补丁最终被实验判定为正确，而不是仅仅自洽。

\step{一句话模型}

\begin{oneline}
电荷是电场唯一的源头，磁场没有源头；变化的磁场卷起电场，电流和变化的电场卷起磁场——电与磁是同一张“电磁场”的两个侧面，而这张场自己携带能量和动量，可以脱离电荷独立存在、独立远行。
\end{oneline}

这句话值得读三遍。第一遍读“源头”：四条方程里有两条专管“源”，电荷独占了电场的源头，磁场的源头一栏至今是空的。第二遍读“卷起”：电场和磁场的旋涡由对方的时间变化和电流共同决定，这是时间变化和空间变化的联立约束，不是接力赛。第三遍读“独立”：场一旦存在，就有自己的能量、动量和演化规律，电荷只是它的源头和受体，不是它的容器。

\step{数学翻译器}

现在把四条总账写成式子。每一条都回答四个问题：描述什么、为什么这样写、何时成立、何时失效；写完立刻用特殊情况检查。

\textbf{方程一（电场的高斯定律）：}
\[
\oint_S \vec E\cdot d\vec A=\frac{Q_{\text{内}}}{\varepsilon_0}
\]
左边读作“闭合面 \(S\) 上电场的净流出”：把面上每一小片处的电场沿外法向的分量累加起来。右边是面内净电荷除以 \(\varepsilon_0\)。为什么这样写：平方反比律配上叠加原理，使得“穿出量”只由总电荷决定，与电荷在面内怎么摆放无关（数学桥层给出理由）。何时成立：在麦克斯韦框架内永远成立，无论场随不随时间变化。何时失效：作为经典定律本身没有失效案例；但“电场线根数正比于电荷”的图示只是示意图，别拿它去数线。检查：\(Q_{\text{内}}=0\) 时净穿出为零——注意是“净”，面内完全可以有电场线穿进穿出；电荷在面外时，它贡献的每一条电场线穿进又穿出，净值为零；把面内电荷整体反号，穿出量反号，方向检查通过。

\textbf{方程二（磁场的高斯定律）：}
\[
\oint_S \vec B\cdot d\vec A=0
\]
描述什么：磁场的净流出恒为零，即磁场无源。为什么右边是零：因为所有寻找孤立磁荷（磁单极子）的实验至今一无所获——这是一个实验事实，不是逻辑必然。何时成立：所有已知实验。何时失效：如果明天发现了磁荷，右边就要改成“面内磁荷”除以某个常数——方程的结构乐于接受这个修改，是实验一直说“不”。检查：把整块磁铁放进面内，从 N 极穿出的磁感线全部绕回 S 极或穿出面外再绕回，净值仍是零；把磁铁换成通电线圈，结论不变。

\textbf{方程三（法拉第定律）：}
\[
\oint_L \vec E\cdot d\vec l=-\frac{d\varPhi_B}{dt},\qquad \varPhi_B=\int_S \vec B\cdot d\vec A
\]
左边是“沿回路 \(L\) 推着单位电荷走一圈，电场做的总功”，即电动势；右边是回路所围曲面上磁通量变化率的负值。第 23 章已详细讲过它描述什么、怎样用，这里补两个边界问题。何时成立：回路固定时直接成立。何时要小心：回路本身在磁场里运动时，总电动势里还有“动生”部分（本质是洛伦兹力，第 22 章），计算时别把两部分重复计入，也别漏掉。检查：\(\varPhi_B\) 不变时环流为零——静电场沿任意回路走一圈做功为零，这正是第 20 章能定义“电势”的前提，新旧知识自洽；磁场均匀增强时得到恒定电动势，电磁炉就是利用这一点在锅底卷起持续的涡旋电场；磁场方向反转，电动势反号，楞次定律的负号检查通过。

\textbf{方程四（安培—麦克斯韦定律）：}
\[
\oint_L \vec B\cdot d\vec l=\mu_0\left(I_{\text{穿}}+\varepsilon_0\frac{d\varPhi_E}{dt}\right)
\]
描述什么：磁场的旋涡由两类“源”维持——穿过曲面的传导电流，和穿过曲面的电通量变化率。为什么这样写：少了第二项，电容器那两个面就给出矛盾答案；添上它，电荷守恒和磁场总账同时闭合。何时成立：经典范围内普适。何时失效：见“模型的边界”。检查：场不随时间变化时，\(\varepsilon_0\,d\varPhi_E/dt=0\)，方程退回旧的安培环路定律——旧定律不是被推翻，而是作为特例被收编；对充电电容器分别取平面 A 和碗面 B，前面已算过，两边都等于 \(\mu_0 I\)，矛盾消除；在导线\textbf{内部}，电通量变化极慢、\(\varepsilon_0\) 又极小（约 \(8.85\times10^{-12}\)），位移电流项与传导电流相比小得可以忽略，所以低频电路里旧安培定律一直很好用——这解释了为什么矛盾等了多年才被人认真对待。

接下来是两条关于“场本身”的公式。

\textbf{场的能量密度：}
\[
u=\frac{1}{2}\varepsilon_0 E^{2}+\frac{B^{2}}{2\mu_0}
\]
描述什么：单位体积的电磁场储存多少能量。为什么这样写：第一项正是第 20 章电容器储能折算到单位体积的结果，第二项正是第 23 章电感储能折算到单位体积的结果——不是新假设，是两笔旧账合并。检查：没有场的地方 \(u=0\)；场强加倍，能量变四倍——平方意味着无论电场朝哪个方向，存的能量都是正的，方向反转检查通过。

\textbf{坡印廷向量（能流密度）：}
\[
\vec S=\frac{1}{\mu_0}\,\vec E\times\vec B
\]
描述什么：单位时间、垂直穿过单位面积的电磁能量，单位是瓦每平方米；对任意闭合面把 \(\vec S\) 累加起来，就是每秒流出这个面的电磁能。为什么用叉乘：只有电场和磁场同时存在、且方向不平行的地方，才有净能流——纯电场或纯磁场只储能，不运输。何时要小心：严格说，给 \(\vec S\) 加上任意一个“原地打转”的场，所有闭合面的总账都不变，所以“能量在空间里走的确切路线”只有总账部分是物理的；但方向和大小的图像在工程上好用到不可替代。检查：\(\vec E\) 与 \(\vec B\) 平行时 \(\vec S=0\)，能量原地不动；\(\vec B\) 反向，\(\vec S\) 反向，运输方向跟着掉头，方向检查通过。

来一组真实数量级。地球轨道处阳光的能流密度约为 \(S\approx1360\) 瓦每平方米。对电磁波（第 25 章给出电场与磁场的配比关系）折算，相当于电场峰值约 \(1000\) 伏每米——和家用电线附近的电场同量级；磁场峰值约 \(3.3\) 微特斯拉，不到地磁场的十分之一。我们每天泡在阳光的电场里而毫无感觉，因为它的方向每秒翻转约 \(10^{15}\) 次，而人体整体是电中性的。

\begin{mathbridge}
先说位移电流的严格推导。对“碗面”B 用电场的高斯定律：穿过 B 的电通量等于碗内净电荷除以 \(\varepsilon_0\)。碗罩住的是上极板，导线每秒往碗里送 \(I\) 库仑，又没有电荷从碗底漏走，所以碗内电荷以速率 \(I\) 增长，于是 \(\varepsilon_0\,d\varPhi_E/dt=I\)。位移电流不是猜出来的对称项，是高斯定律加上电荷守恒逼出来的。

再把积分语言翻译成微分语言。取一个闭合面，把它越缩越小，缩到一个点：净流出除以体积的极限，叫该点的\textbf{散度}；取一个回路缩到一个点：环流除以面积的极限（按方向取最大者），叫该点的\textbf{旋度}。用散度定理和旋度定理，可以把四条积分方程逐点改写。电场高斯定律变成“电场的散度等于电荷密度除以 \(\varepsilon_0\)”；法拉第定律变成“电场的旋度等于磁场时间变化率的负值”。另外两条的翻译放进挑战层。
\end{mathbridge}

\begin{challenge}
逐点形式的麦克斯韦方程组：
\[
\nabla\cdot\vec E=\frac{\rho}{\varepsilon_0},\qquad
\nabla\cdot\vec B=0,
\]
\[
\nabla\times\vec E=-\frac{\partial\vec B}{\partial t},\qquad
\nabla\times\vec B=\mu_0\vec J+\mu_0\varepsilon_0\frac{\partial\vec E}{\partial t}.
\]
请盯着两个旋度方程看：它们把 \(\vec E\) 和 \(\vec B\) 的\textbf{空间}变化与\textbf{时间}变化交叉锁死。在没有任何电荷电流的区域（\(\rho=0\)，\(\vec J=0\)），对第三条两边取旋度、代入第四条（或反过来做），配一条矢量恒等式，会得到 \(\vec E\) 和 \(\vec B\) 各自满足波动方程，波速为
\[
v=\frac{1}{\sqrt{\mu_0\varepsilon_0}}.
\]
代入当时的测量值，\(v\approx3\times10^{8}\) 米每秒——与已知光速在误差内吻合。麦克斯韦 1865 年由此断言光就是电磁波。注意整个推导里没有出现“电场生磁场、磁场生电场”的接力图像：波动解是\textbf{同时}满足全部四条约束的时空分布。

还有一道加餐：假如磁荷存在，密度记为 \(\rho_m\)，那么第二式右边变成正比于 \(\rho_m\)，第三式要添一项正比于磁流的项，方程组会变得在电、磁之间几乎完全对称。狄拉克后来还证明，只要宇宙里存在一个磁单极子，电荷的量子化就能得到解释。这个诱人的对称至今没有实验证据——方程二右边的零，仍然是一笔等待宣判的账。
\end{challenge}

\step{模型的边界}

第一，位移电流不是“电荷在真空里流动”。极板之间没有载流子，没有漂移速度，不产生焦耳热；它只是在“产生磁场”和“闭合电荷守恒的账”这两件事上扮演了电流的角色。名字里的“位移”来自麦克斯韦当年设想的介质弹性机械模型——那个模型早已被放弃，名字留了下来。把它\textbf{当成}什么：一种在磁效应上与电流等效的变化率。别把它当成什么：别当成有粒子在挪动，也别指望用电流表在真空中测到它。

第二，电路定律有适用边界。基尔霍夫定律假设“同一瞬间，一条支路上的电流处处相等”，这只在\textbf{似稳}条件下成立：场扰动传遍整个电路所需的时间，要远小于电流变化的特征时间。电磁波一秒能绕地球七圈半，对桌面电路，传遍一遍只要几纳秒，远快于市电 \(50\) 赫兹的变化，似稳近似好得惊人；但对几千公里的输电线路、吉赫兹的手机射频电路，必须回到完整的麦克斯韦方程——工程上叫传输线理论。判据一句话：系统尺寸要远小于“扰动速度乘以特征时间”。

第三，麦克斯韦方程只管场，不管料。电荷在场里怎么动，要另配洛伦兹力（第 22 章）；材料里的电荷怎么响应场，要另配介电、导电的规律。方程也不回答“电荷为什么是一份一份的”“\(\varepsilon_0\) 和 \(\mu_0\) 为什么是这个值”。它是一张极其严格的网，网住的是场与源的关系，不是全部物理。

第四，经典场本身有失效处。光极弱时，能量是一份一份到达探测器的；光电效应和黑体辐射的灾难（第 42 章）最终要求把电磁场本身量子化。在极强场下，真空还会出现经典方程无法描述的新现象。但在从收音机到电网的全部日常与工程范围里，麦克斯韦方程精确得近乎无情——它是经典物理里活得最好的理论之一。

第五，以上四条方程的写法假定真空。到了介质内部，\(\varepsilon_0\)、\(\mu_0\) 要被材料自己的参数替换，电磁扰动的传播速度也随之变慢——这埋下了“光在玻璃里为什么走得慢”的伏笔，第 36 章会捡起它。

\step{回头解释开场现象}

现在回到台灯。用坡印廷向量把能量走的路完整走一遍。

电池把正负电荷分别泵到电路两端，于是整个导线网络的表面上分布着一层微妙的面电荷。这层电荷在导线内外建立起电场：导线内部，\(\vec E\) 沿导线方向，驱动电流（第 21 章）；紧靠导线表面的地方，\(\vec E\) 也大致沿导线。电流的磁场 \(\vec B\) 环绕导线。于是 \(\vec S=\vec E\times\vec B/\mu_0\)：在电池附近，\(\vec S\) 指向电池\textbf{外}——能量从电池的“侧壁”涌入空间；在连接导线旁，\(\vec S\) 大致沿导线指向灯泡——能量在导线\textbf{外面}的空间里被护送前进；在灯丝旁，\(\vec S\) 指向灯丝\textbf{内部}——能量从侧面涌入，变成热和光。

\begin{center}
\begin{tikzpicture}[>=Stealth]
  \draw[thick] (0,0) -- (6.4,0);
  \draw[->,thick] (0.5,0.35) -- (1.7,0.35) node[midway,above]{电流 $I$};
  \draw[->,thick,blue!60!black] (2.5,0.95) -- (3.9,0.95) node[midway,above]{$\vec E$ 沿导线};
  \draw[thick] (3.2,0) ellipse (0.26 and 0.72);
  \node at (3.2,-1.15) {$\vec B$ 环绕导线};
  \draw[->,very thick,red!70!black] (5.1,1.25) -- (4.4,0.22);
  \draw[->,very thick,red!70!black] (5.1,-1.25) -- (4.4,-0.22);
  \node[red!70!black] at (5.6,0.75) {$\vec S$};
  \node[red!70!black] at (5.35,-0.85) {能量从侧面流入};
\end{tikzpicture}
\end{center}

这笔账可以验算。以一盏 \(60\) 瓦的台灯为例：电压 \(220\) 伏，电流约 \(0.27\) 安。在距离电源线 \(1\) 毫米处，电流的磁场约为 \(5\times10^{-5}\) 特斯拉，只有地磁场的千分之一；而“护送”能量的电场也就几百伏每米的量级。看不见、摸不着，但它们每秒向灯丝方向输送的，正是实实在在的 \(60\) 焦耳。

再做一个更干净的理论验算。设一段长 \(L\) 的导线，两端电压降为 \(V\)，电流为 \(I\)，导线半径为 \(a\)。表面处电场 \(E=V/L\)，磁场 \(B=\mu_0 I/(2\pi a)\)，于是能流密度
\[
S=\frac{EB}{\mu_0}=\frac{VI}{2\pi aL},
\]
方向指向导线内部。乘以导线的侧面积 \(2\pi aL\)，每秒流入这段导线的总能量恰好是 \(VI\)——正好等于这段导线消耗的焦耳热功率，一分不差。这是坡印廷图像最硬的一次特例检查：从空间流进的，和导线内部热掉的，账面严丝合缝。

再交代“铁轨”比喻的边界。它\textbf{像的部分}：导线决定能量去哪里，改线就换目的地，正如铁轨决定火车开往哪里。它\textbf{不像的部分}：火车开在铁轨上面，能量却走在导线\textbf{两侧的空间里}；更要命的是，导线里电子的漂移速度只有每秒零点几毫米，而能量以接近光速的速度沿空间送达——“搬运工”图像在这两个数字面前彻底出局。你在“先猜一猜”里若选了（c），现在可以对上号了。

太阳和收音机也一并结账。真空中没有电荷、没有电流，麦克斯韦方程依然允许电场和磁场结伴的波动解，以 \(1/\sqrt{\mu_0\varepsilon_0}\) 的速度前进，不需要任何介质——阳光就是这样走完一亿五千万公里的。动手实验里的“咔哒”声也一样：开关和火花让电流在极短时间里突变，突变的电流与电场向空间发出一小团电磁扰动，中波收音机的天线把它捡起来、放大成声音。至于这团扰动为什么偏偏以光速前进、它和光凭什么是一家人——那是第 25 章的故事。

\step{脑洞实验与挑战题}

\begin{thought}[追不上的光]
麦克斯韦方程刚出世不久，十六岁的爱因斯坦给自己出了一道题：假如我能以光速追着一束电磁波跑，我会看见什么？

按直觉，追平之后，波应该“静止”下来：我会看到一个不随时间变化、只在空间上波浪起伏的电场和磁场分布——像一串凝固的涟漪。

现在把麦克斯韦方程当裁判。不随时间变化的场，方程三、四右边的时间变化率全为零，于是电场、磁场都不允许有旋涡——“凝固的涟漪”这种分布根本不满足方程，是不合法的解。要么方程错了，要么“以光速追平一束光”这个前提本身有问题。结局写在第 38 章：方程没错，是时间和空间的概念被改了。麦克斯韦方程组因此不只是电磁学的总结，它还是通往相对论的大门。
\end{thought}

电磁场既然携带能量，也携带动量，光照在物体上会产生压力——光压。地球轨道处，阳光照在理想反射面上的压强约为 \(9\) 微帕，小得像开玩笑。但把尺度推到极端就有戏看了：一面 \(14\) 米见方的太阳帆，承受的光压推力约 \(2\) 毫牛，只相当于几粒芝麻的重量；可它不用燃料、日夜不停地推，日本的 IKAROS 探测器就靠它在 2010 年完成了行星际航行。再算笔激光的账：\(1\) 千瓦的激光照在全反射镜上，推力 \(F=2P/c\approx7\) 微牛；想用光压举起一枚鸡蛋（约 \(0.5\) 牛），需要约 \(75\) 兆瓦的激光——一座大型电站的输出功率，全变成光。场有动量，在麦克斯韦的时代是纸上结论，在今天的航天工程里是设计参数。

\begin{puzzles}
\item 关掉台灯的那一瞬间，原来导线周围空间里的电磁场能量去哪了？提示：场能不能“瞬间”消失？扰动以有限速度传开，大部分能量在熄灭过程中流进灯丝变成最后的热，还有一小部分以电磁波的形式被带走——所以开关火花能在收音机里听到。
\item 充电中的电容器，极板间隙里的位移电流卷出的磁场，绕向与导线中传导电流卷出的磁场一致还是相反？若把极板面积加大、充电电流保持不变，间隙里同一位置的磁场会怎样变？提示：先用右手定则比方向，再想“电流的账必须连续”——面积加大后位移电流摊开变薄，局部磁场变弱，但以回路为边界的总账仍等于 \(\mu_0 I\)。
\item 一只充好电的大电容器静静放在桌上，极板间有强大但恒定不变的电场，储存着可观的能量。它周围会有磁场吗？它会持续向外辐射能量吗？提示：查方程四，\(d\varPhi_E/dt\) 是多少；再查坡印廷向量，\(\vec B\) 为零时 \(\vec S\) 是多少。储能不等于辐射。
\item 假如明天的实验确凿发现了磁荷，四条方程哪些要改、怎么改？改完之后，电和磁在方程里的地位是更对称还是更不对称？提示：先想“磁荷是磁场的源”该写进哪一条；再想“运动的磁荷”（磁流）会不会像电流卷出磁场那样卷出电场，又该写进哪一条。
\end{puzzles}

\section*{本章小结}

麦克斯韦没有做任何新实验。他把库仑、安培、法拉第的定律摆在一起，发现它们在电容器面前自相矛盾；为了补上裂缝，他让“变化的电场”在产生磁场这件事上与电流平起平坐——位移电流由此而来。补好之后的四条方程告诉我们：电荷是电场唯一的源头，磁场没有源头，变化的磁场卷起电场，电流和变化的电场卷起磁场。电磁场携带能量与动量，可以脱离电荷在真空中远行，能量走的是空间而不是导线。这四行方程是十九世纪物理学的最高成就，也是两份礼物的包装纸：一份叫电磁波（第 25 章），一份叫相对论（第 38 章）。
